\chapter[Minihalos]{Minihalos}
\label{Chap:Background}

Since the first discoveries of diffuse radio emission, much research has been devoted towards understanding the mechanisms behind the production of synchrotron electrons. These electrons are also known as `cosmic ray' (CR) electrons or non-thermal electrons to differentiate them from the thermal electrons of the ICM. In the case of CR electrons within the thin plasma of the ICM, we make the standard assumption that the population of CR electrons is homogeneous and isotropic and follows a power-law energy distribution $n(E) \propto E^{-\delta}$. From this we then predict the following basic observables (see \citet{vanWeeren2019}, \citet{Ferreti2012} and references within). \stickynote{Add graph of spectral shapes?}

\begin{enumerate}
    \item The synchrotron emissivity is dependent on the number density of the CR electrons and the strength of the magnetic field. As such, the emissivity also follows a power-law $S(\nu) \propto \nu^{-\alpha}$ which is related to $\delta$ by $\alpha = (\delta - 1)/2$. Typical minihalo spectral indices are steep with $\alpha > 1$.
    \item Over time the emitting particles lose energy, which changes the particle energy distribution and therefore, the emitted radio spectrum. This appears as a steepening or curved spectrum, rather than a flat power-law. The break frequency at which this occurs is related to the time of acceleration. This means that radio events that have stopped will show an aged spectrum with a characteristic break, such as radio relics, while ongoing events are continually injecting new particles and don't feature the same steepening due to age. 
    \item Radiation from CR electrons passing through a uniform magnetic field is linearly polarised. Non-uniform and disordered fields reduce the linear polarisation. 
\end{enumerate} 

\section{Production mechanisms}

The 

\subsection{Turbulent re-acceleration}

\label{sec:BG_reacceleration}
\begin{itemize}
    \item What is the process here.
    \item How manifest spectrally.
    \item How manifest spatially.
\end{itemize}

\subsection{Hadronic Production}
\label{sec:BG_hadronic}
\begin{itemize}
    \item What is the process for hadronic production
    \item How does this manifest spectrally and evolve. What should we expect observationally here based on simulation work.
    \item How does this manifest spatially and evolve, what should we see observationally.
\end{itemize}



\subsection{Large Radio Halos}
\begin{itemize}
    \item Use similar physics but much larger scale. 
    \item Some basics things about previous research in differentiating production. 
    \item Gamma ray non-detection.
    \item General consensus is that its re-acceleration.
\end{itemize}

\subsection{Minihalos}
\begin{itemize}
    \item These are much smaller, so therefore much harder to observe. Difference to large radio halos.
    \item Similarities between host clusters (CC), general properties of spectral index etc.
\end{itemize}









\subsection{Previous minihalo studies}
\begin{itemize}
    \item Generally look at lower frequencies, or a singular study for a particular minihalo.
    \item Want to look at more minihalos in a more holistic manner. 
    \item Some findings from previous papers e.g. RXJ1720, MS1455+2232 in basic detail. (more detail on the per-halo discussion).
\end{itemize}
